\documentclass[12pt]{article}
\usepackage[utf8]{inputenc}
\usepackage[T1]{fontenc}
\usepackage[spanish]{babel}
\usepackage{amsmath,amssymb}
\usepackage{geometry}
\usepackage{multicol}
\usepackage{enumitem}
\usepackage{xcolor}
\usepackage{lmodern}
\geometry{letterpaper, margin=2cm}

\title{\textbf{Solución Guía 02}\\ Potencias y Notación Científica}
\author{}
\date{}

\begin{document}
	
	\maketitle
	
	\section*{Parte I: Simplificación de potencias}
	
	\begin{enumerate}
		
		\item $
		\dfrac{2^5\cdot 2^3\cdot 2^{-2}}{2^4}
		=2^{5+3-2-4}
		=2^2
		=4
		$
		
		\item $
		(3^2\cdot 3^4)^2
		=(3^{2+4})^2
		=(3^6)^2
		=3^{12}
		=531441
		$
		
		\item $
		\dfrac{5^7}{5^3\cdot 5^{-1}}
		=\dfrac{5^7}{5^{3+(-1)}}
		=\dfrac{5^7}{5^2}
		=5^{7-2}
		=5^5
		=3125
		$
		
		\item $
		\dfrac{(2^3)^4}{2^5}
		=\dfrac{2^{12}}{2^5}
		=2^{12-5}
		=2^7
		=128
		$
		
		\item $
		\left(\dfrac{4^3\cdot 4^2}{4^4}\right)^2
		\left(\dfrac{2^5}{2^2}\right)^3
		=
		\left(4^{3+2-4}\right)^2
		\left(2^{5-2}\right)^3
		=
		(4^1)^2(2^3)^3
		=4^2\cdot 2^9
		=16\cdot 512
		=8192
		$
		
	\end{enumerate}
	
	\section*{Parte II: Operaciones en notación científica}
	
	\begin{enumerate}
		
		\item $
		(4\times10^5)(3\times10^3)
		=(4\cdot 3)\times10^{5+3}
		=12\times10^8
		=1.2\times10^9
		$
		
		\item $
		\dfrac{6\times10^8}{2\times10^4}
		=\left(\dfrac{6}{2}\right)\times10^{8-4}
		=3\times10^4
		$
		
		\item $
		(2.5\times10^3)(4\times10^{-2})
		=(2.5\cdot 4)\times10^{3+(-2)}
		=10\times10^1
		=1\times10^2
		$
		
		\item $
		\dfrac{9\times10^{-4}}{3\times10^{-2}}
		=\left(\dfrac{9}{3}\right)\times10^{-4-(-2)}
		=3\times10^{-2}
		$
		
	\end{enumerate}
	
	\section*{Parte III: Exprese en notación científica las siguientes cantidades}
	
	\begin{enumerate}
		
		\item $
		4.500.000.000=4.5\times10^9
		$
		
		\item $
		8.200.000=8.2\times10^6
		$
		
		\item $
		35.000.000=3.5\times10^7
		$
		
		\item $
		7.600.000=7.6\times10^6
		$
		
		\item $
		950.000=9.5\times10^5
		$
		
		\item $
		0.00045=4.5\times10^{-4}
		$
		
		\item $
		0.0000028=2.8\times10^{-6}
		$
		
		\item $
		0.00003=3\times10^{-5}
		$
		
		\item $
		0.0000000015=1.5\times10^{-9}
		$
		
		\item $
		0.00000032=3.2\times10^{-7}
		$
		
		\item $
		0.0007=7\times10^{-4}
		$
		
		\item $
		0.00000008=8\times10^{-8}
		$
		
	\end{enumerate}
	
	\section*{Parte IV: Aplicaciones en Agronomía}
	
	\textbf{Nota:} En el archivo original hay un ejercicio sin \verb|\item|. Aquí se presenta la versión corregida y resuelta.
	
	\begin{enumerate}[start=9]
		
		\item En una parcela experimental se plantan \(3^5\) semillas por sector.  
		Si el terreno tiene \(3^3\) sectores iguales, determine el número total de semillas sembradas.
		
		\[
		3^5\cdot 3^3=3^{5+3}=3^8=6561
		\]
		
		\textbf{Respuesta:} Se siembran \(\boxed{6561}\) semillas.
		
		\item Una colonia de bacterias del suelo se duplica cada hora.  
		Si inicialmente hay \(2^4\) bacterias, ¿cuántas bacterias habrá después de 8 horas?
		
		\[
		2^4\cdot 2^8=2^{12}=4096
		\]
		
		\textbf{Respuesta:} Habrá \(\boxed{4096}\) bacterias.
		
		\item Un fertilizante contiene \(4\times10^{-3}\) kg de micronutriente por kilogramo de producto.  
		Si se aplican \(3\times10^3\) kg de fertilizante, determine la cantidad total de micronutriente aplicada.
		
		\[
		(4\times10^{-3})(3\times10^3)
		=12\times10^0
		=12
		\]
		
		\textbf{Respuesta:} Se aplican \(\boxed{12\text{ kg}}\) de micronutriente.
		
		\item Un fertilizante contiene \(8\times10^{-4}\) kg de micronutriente por kilogramo de producto.  
		Si se aplican \(2\times10^4\) kg de fertilizante, determine la cantidad total de micronutriente aplicada.
		
		\[
		(8\times10^{-4})(2\times10^4)
		=16\times10^0
		=16
		\]
		
		\textbf{Respuesta:} Se aplican \(\boxed{16\text{ kg}}\) de micronutriente.
		
		\item Un cultivo utiliza \(5^4\) semillas por parcela.  
		Si existen \(5^3\) parcelas en el terreno, ¿cuántas semillas se utilizan en total?
		
		\[
		5^4\cdot 5^3=5^{4+3}=5^7=78125
		\]
		
		\textbf{Respuesta:} Se utilizan \(\boxed{78125}\) semillas.
		
	\end{enumerate}
	
\end{document}